\begin{figure}[ht!]
\centering
\begin{tikzpicture}
\tikzset{
    grow'=right,
    level distance=48pt,
    sibling distance=6pt,
    execute at begin node=\strut,
    every tree node/.style={anchor=base west},
    edge from parent/.append style={very thick},
    edge from parent/.style={
        draw,
        edge from parent path={
            (\tikzparentnode.east) -| ($(\tikzparentnode.east)!0.5!(\tikzchildnode.west)$) |- (\tikzchildnode.west)
        },
    },
    every node/.style={anchor=center,font=\small\bfseries, text centered, inner sep=0pt},
    every level 0 node/.style={circle, font=\small\bfseries, draw, fill=blue!30, inner sep=0pt},
    every internal node/.style={font=\small, inner sep=4pt},
    every leaf node/.style={rectangle, draw, fill=blue!30, minimum width=2.5cm, text centered},
    frontier/.style={distance from root=\textwidth-20pt-100pt},
}
\Tree [.\(I\)
    [.\(F_1^{\text{succ}}\)
        [.\(F_2^{\text{succ}}\)
            [.\(X_1\) ]
        ]
        [.\(F_2^{\text{fail}}\)
            [.\(X_2\) ]
        ]
    ]
    [.\(F_1^{\text{fail}}\)
        [.\(X_3\) ]
    ]
]
\end{tikzpicture}
\caption{Illustrative event tree with an initiating event \(I\), two functional events \(F_1\) and \(F_2\), and three end-states \(X_1, X_2, X_3\).}
\label{fig:event_tree_example}
\end{figure}